\documentclass[11pt]{article}
\usepackage[T1]{fontenc}
\usepackage[margin=1in]{geometry}
\usepackage{amsmath,amssymb,amsthm}
\usepackage{microtype}
\usepackage[hidelinks]{hyperref}
\setlength{\parindent}{0pt}
\setlength{\parskip}{0.3em}
\newtheorem{theorem}{Theorem}
\newtheorem{lemma}[theorem]{Lemma}
\newtheorem{corollary}[theorem]{Corollary}
\newcommand{\R}{\mathbb R}
\newcommand{\vol}{\lambda}
\newcommand{\Int}{\operatorname{int}}
\newcommand{\code}[1]{\texttt{\detokenize{#1}}}
\title{Disproof of Conjecture 00000001624\\[0.3em]
\large A positive-measure Cantor spectrum must have an internal gap}
\author{C0ldSmi1e\\\small AI-assisted submission}
\date{9 October 2026}

\begin{document}
\maketitle
\begin{abstract}
The conjecture requests a quasiperiodic potential whose spectrum is a
positive-measure Cantor set with no gaps and whose spectral measure is singular.
Its geometric requirements are already incompatible under the literal,
standard real-energy meaning of ``with no gaps.'' We prove that every closed
subset of $\R$ with empty interior and positive Lebesgue measure has a bounded
complementary component with both endpoints in the set. The Lean proof ranges
over all real sets and then over every potential type and spectrum assignment;
no particular operator or potential is selected.
\end{abstract}

\section{Exact claim and interpretation}

The English source states:
\begin{quote}
\raggedright
``Conjecture: There exists a quasiperiodic potential whose spectrum is a
positive-measure Cantor set with no gaps while the spectral measure is singular
(an explicit construction of gapless singular spectra).''
\end{quote}
The accompanying Chinese statement also attributes both Cantor character and
absence of gaps to the spectrum. The exact original in both languages is
retained in \code{source/ORIGINAL.md}. Neither version specifies a reference
energy at which alone a gap should be absent.

Throughout this report the spectrum is a subset of the real energy line, as
in the usual self-adjoint setting for quasiperiodic potentials. A real Cantor
set means a nonempty, compact, perfect, nowhere-dense subset of $\R$. It need
not be the middle-thirds Cantor set: Cantor character does not by itself imply
zero Lebesgue measure. The measure in ``positive-measure'' is Lebesgue measure
$\vol$ on $\R$, distinct from the operator's spectral measure.

For $S\subseteq\R$, an \emph{internal gap} is a nonempty interval $(a,b)$ such
that
\begin{equation}\label{eq:gap}
 a<b,\qquad a,b\in S,\qquad (a,b)\cap S=\varnothing.
\end{equation}
These finite spectral endpoints distinguish an internal gap from the exterior
resolvent rays of a bounded spectrum. ``With no gaps'' means that no such pair
$a,b$ exists, anywhere in the spectrum. We verify below that
\eqref{eq:gap} yields a whole connected component of the complement. For closed
real sets, we also prove that this absence of gaps is equivalent to containing
every point between any two of the set's points.

This reading differs from absence of a gap at a chosen Fermi energy, absence
of an excitation gap above a selected ground state, or closure of only some
labelled gaps. No such restriction is present in the source. The result here
does not settle a rewritten conjecture with one of those different conditions.

\section{The set-level obstruction}

\begin{lemma}[A missing point lies in a full internal gap]\label{lem:gap}
Let $S\subseteq\R$ be closed, and let $x,y\in S$ and $c\notin S$ satisfy
$x<c<y$. Then there are $a,b\in S$ with $a<c<b$ and
$(a,b)\cap S=\varnothing$.
\end{lemma}
\begin{proof}
Set
\[
 L=S\cap(-\infty,c],\qquad R=S\cap[c,\infty),\qquad
 a=\sup L,\qquad b=\inf R.
\]
The sets $L,R$ are nonempty because they contain $x,y$, respectively.
They are closed; $L$ is bounded above by $c$, and $R$ is bounded below by $c$.
The supremum and infimum therefore belong to their respective sets.
Consequently $a,b\in S$ and $a\leq c\leq b$. Neither equality can hold because
$c\notin S$, so $a<c<b$.

If $z\in S\cap(a,b)$ and $z\leq c$, then $z\in L$ and hence $z\leq a$,
a contradiction. If instead $c\leq z$, then $z\in R$ and $b\leq z$, also a
contradiction. Thus $(a,b)$ is disjoint from $S$.
\end{proof}

\begin{lemma}[The interval is an entire complementary component]\label{lem:component}
If \eqref{eq:gap} holds and $c\in(a,b)$, then the connected component of
$\R\setminus S$ containing $c$ is precisely $(a,b)$.
\end{lemma}
\begin{proof}
The interval $(a,b)$ is connected, contains $c$, and lies in $\R\setminus S$,
so it lies in that component. Conversely, a connected subset of $\R$ contains
every point between any two of its points. If the component contained some
$z\leq a$, it would contain $a$ between $z$ and $c$, contradicting $a\in S$.
If it contained some $z\geq b$, it would similarly contain $b\in S$.
The component is therefore contained in $(a,b)$.
\end{proof}

\begin{theorem}\label{thm:main}
Every closed subset $S\subseteq\R$ with empty interior and
$\vol(S)>0$ has an internal gap.
\end{theorem}
\begin{proof}
A set with at most one point has Lebesgue measure zero. Thus $S$ contains
distinct points, which we order as $x<y$. The nonempty open interval $(x,y)$
cannot be contained in $S$: otherwise it would be contained in $\Int(S)$,
contrary to $\Int(S)=\varnothing$. Choose $c\in(x,y)\setminus S$ and apply
Lemma~\ref{lem:gap}. Lemma~\ref{lem:component} confirms that the resulting
gap is a bounded complementary component, not an arbitrarily selected
subinterval of a larger one.
\end{proof}

The argument actually needs only two distinct points in place of positive
measure, and does not require compactness or perfection. Those observations
strengthen the obstruction; they do not add assumptions to the conjecture.

\begin{lemma}[Equivalent no-gap formulation]\label{lem:convex}
For a closed set $S\subseteq\R$, absence of internal gaps is equivalent to
order convexity: whenever $x,y\in S$ and $x\leq z\leq y$, one has $z\in S$.
\end{lemma}
\begin{proof}
If $z\notin S$ lies between $x,y\in S$, then it differs from both endpoints,
so $x<z<y$; Lemma~\ref{lem:gap} supplies a gap. Conversely, order convexity
would put every point of an interval $(a,b)$ with $a,b\in S$ into $S$,
contradicting \eqref{eq:gap}.
\end{proof}

\section{Negation of the original existence claim}

\begin{corollary}\label{cor:cantor}
There is no positive-measure real Cantor set with no internal gaps.
\end{corollary}
\begin{proof}
A real Cantor set is closed and has empty interior. Apply
Theorem~\ref{thm:main}.
\end{proof}

Every proposed witness to the source claim would provide such a real spectrum,
contradicting Corollary~\ref{cor:cantor}. This is a disproof of a necessary
conjunction of the original conditions, hence of the full existential claim.
The additional quasiperiodicity and spectral-measure singularity requirements
cannot restore a witness once its required spectrum is impossible. There can
therefore be no explicit construction meeting all the stated requirements.
No separate claim about the existence of singular spectral measures is made.

The formal lift makes the quantifier coverage explicit. Let $\mathcal V$ be
\emph{any} type of candidate potentials, let
$\Sigma:\mathcal V\to\mathcal P(\R)$ be \emph{any} real spectrum assignment
(where the codomain denotes the power set), and let $Q,T$ be arbitrary
predicates on $\mathcal V$. Then
\[
\neg\exists V\in\mathcal V:\quad
Q(V)\ \land\ \operatorname{Cantor}(\Sigma(V))\ \land\
\vol(\Sigma(V))>0\ \land\ \operatorname{NoGaps}(\Sigma(V))\ \land\ T(V).
\]
Here genuine quasiperiodicity and singular-spectral-measure predicates may be
substituted for $Q,T$. Universality of this result means that these notions
are not replaced by a restricted toy model. The proof works even without
constraining $Q$ or $T$ at all. The source's omitted choices of operator
realization, potential regularity, phase, and spectral-measure vector do not
affect this geometric contradiction.

\section{Correspondence with the Lean project}

The project uses Lean 4.19.0 and Mathlib v4.19.0. Its mathematical core is
\code{Conjecture1624.lean}, in namespace \code{Conjecture1624}.
\begin{itemize}
\setlength{\itemsep}{0.2em}
\item \code{IsRealCantor} is nonemptiness, compactness, Mathlib's
\code{Perfect}, and empty interior. Mathlib's \code{Perfect} includes
closedness, so the last condition is equivalent here to nowhere density.
\item \code{IsGap} is exactly \eqref{eq:gap}; \code{HasNoGaps} is its
existential negation.
\item \code{exists_gap_around} proves Lemma~\ref{lem:gap},
and \code{IsGap.component} proves Lemma~\ref{lem:component}.
\item \code{exists_gap_of_positive_volume} is Theorem~\ref{thm:main}.\\
\code{hasNoGaps_iff_ordConnected} is Lemma~\ref{lem:convex}.
\item \code{no_source_spectrum} is Corollary~\ref{cor:cantor}, and
\code{no_source_potential} is the universal lift just displayed.
\end{itemize}

The conclusion uses ordinary real topology and Lebesgue measure. No numerical
experiment or auxiliary computation is needed to establish the mathematical
claim. Reproducible build, challenge checks, axiom-dependency inspection, and
independent review records accompany the submission. Those checks establish
local verification; repository acceptance remains a maintainer decision.

\end{document}
